\documentclass[10pt,a4paper]{amsart}
\usepackage[margin=19mm]{geometry}
\usepackage{amsmath,amssymb,mathtools}
\usepackage[scr=rsfs,cal=euler]{mathalfa}
\usepackage{xfrac}
\usepackage[unicode,hidelinks,bookmarks=false]{hyperref}
\newcommand{\ie}{i.e.\ }
\newcommand{\cf}{cf.\ }
\newcommand{\resp}{respectively\ }
\newcommand{\Subsection}{\subsection}
\providecommand{\nobreakdash}{\nobreak\hbox{-}}
\setlength{\emergencystretch}{2em}
\multlinegap0pt
\usepackage{dsfont}
\usepackage{amssymb}
\usepackage{xcolor}
\usepackage{tcolorbox}
\newtheorem{proposition}{Proposition}
\title{Structural functional identifiability and model discovery in \\ differential equation models}
\author{Torkel E Loman, Alexander P Browning, Ruth E Baker}
\date{Source: arXiv:2606.30289v1 (CC BY 4.0)}
\begin{document}
\maketitle

\noindent\emph{Source and licence.} Structural functional identifiability and model discovery in \\ differential equation models. Version arXiv:2606.30289v1 (CC BY 4.0), distributed by arXiv under the Creative Commons Attribution 4.0 International licence, http://creativecommons.org/licenses/by/4.0/, as recorded in the arXiv licence metadata for this exact version and shown on the abstract page https://arxiv.org/abs/2606.30289v1. Full licence notice: \texttt{licenses/arxiv-2606.30289-CC-BY-4.0.txt}.\par\smallskip

\noindent\textit{Editorial scope.} Counted unit: the article's proposition prop:diffalg\_one\_parameter\_identifiability (source0.tex lines 270--276). The article's complete original proof of it is delivered with it (source0.tex lines 278--299, one \texttt{proof} environment of 22 lines). No mathematics is written, repaired or re-derived: the statement and the proof are the article's own lines, in the article's own notation, and no retained line is substituted or re-flowed.  Retained uncounted as the source-local context the proof needs: lines 264--267.  Nothing was omitted from any retained span, so the package carries no omission record.  No retained line contains a citation, so the wrapper carries no bibliography and no bibliography entry.  No retained line contains a figure environment, an image-inclusion command or drawing code, so no image is delivered and the package declares no asset dependency.  Editorial formatting only, in addition: an AMS \texttt{amsart} preamble that restates the article's own declarations and macros used by the retained lines, namely the environments \texttt{proposition} on separate counters, together with the standard packages those lines need.  The retained environments print in this document as Proposition 1 in order of appearance, whereas the article prints the same items with the numbering of its own full document.  The complete native source of the article as distributed in the arXiv e-print for this exact version is bundled unchanged as \texttt{sources/204009/source0.tex}.
\begin{align}
    \dot{X}_1 &= f(X_2) - dX_1, \label{eq:diffalg_model1_dXdt} \\
    \dot{X}_2 &= X_1 - dX_2,   \label{eq:diffalg_model1_dYdt}
\end{align}

\begin{proposition} \label{prop:diffalg_one_parameter_identifiability}\label{prop:diffalg_case1}\label{prop:diffalg_case2}
In the single-parameter mutual activation model, Equations~\eqref{eq:diffalg_model1_dXdt}--\eqref{eq:diffalg_model1_dYdt}, both $f$ and $d$ are structurally identifiable in either of the following cases:
\begin{enumerate}
    \item Only $X_2$ is observable.
    \item Only $X_1$ is observable, provided $f$ is locally invertible and $d \neq 0$.
\end{enumerate}
\end{proposition}

\begin{proof}[Proof of part (i)]
To eliminate $X_1$, we first isolate it in Equation~\eqref{eq:diffalg_model1_dYdt} and differentiate:
\begin{equation*}
    X_1= \dot{X}_2 + dX_2 \quad \Longrightarrow \quad \dot{X}_1 = \ddot{X}_2 + d\dot{X}_2.
\end{equation*}
Substituting into Equation~\eqref{eq:diffalg_model1_dXdt} gives the following input-output equation
\begin{equation}
    \label{eq:diffalg_model1_Yobs_dY2dt2}
    \ddot{X}_2 + 2d\dot{X}_2 + d^2 X_2 - f(X_2) = 0.
\end{equation}
Now suppose two pairs $(f_1, d_1)$ and $(f_2, d_2)$ yield identical $X_2$ dynamics. Inserting both into Equation~\eqref{eq:diffalg_model1_Yobs_dY2dt2}, isolating $\ddot{X_2}$, and then equating gives
\begin{equation}
    \label{eq:diffalg_model1_Yobs_equiveq}
    -2d_1\dot{X}_2 - d_1^2 X_2 + f_1(X_2) = \ddot{X}_2 = -2d_2\dot{X}_2 - d_2^2 X_2 + f_2(X_2).
\end{equation}
Here, under the assumption of an arbitrary amount of noiseless data, each of $X_2$'s derivatives can be varied independently, allowing us to separate Equation~\eqref{eq:diffalg_model1_Yobs_equiveq} into two independent equalities that must hold:
\begin{align}
    -2d_1\dot{X}_2       &= -2d_2\dot{X}_2,       \label{eq:diffalg_model1_Yobs_split1} \\
    -d_1^2 X_2 + f_1(X_2) &= -d_2^2 X_2 + f_2(X_2).  \label{eq:diffalg_model1_Yobs_split2}
\end{align}
Equation~\eqref{eq:diffalg_model1_Yobs_split1} immediately gives $d_1 = d_2 \equiv d$. Substituting into Equation~\eqref{eq:diffalg_model1_Yobs_split2} then gives $f_1(X_2) = f_2(X_2)$. This demonstrates that two different pairs $(f, d)$ cannot yield identical dynamics with respect to $X_2$, showing that both $d$ and $f$ are structurally identifiable.
\end{proof}

\end{document}
